\section{Metodologia}
\label{desenvolvimento}

A pesquisa é quantitativa e experimental e segue as fases da CRISP-NET (Figura~\ref{fig-fluxo}). O desenho é fatorial $2 \times 2$: variam só a distância e o método; dados, janelas, $k$, construção do grafo e avaliação são idênticos nas quatro variações (Tabela~\ref{tab-variacoes}).

\begin{figure}[!ht]
\centering
\begin{tikzpicture}[node distance=2mm and 2mm, >=latex,
  box/.style={draw, rounded corners=2pt, align=center, font=\scriptsize, text width=0.84\columnwidth, inner sep=2.2pt},
  meia/.style={draw, rounded corners=2pt, align=center, font=\scriptsize, text width=0.395\columnwidth, inner sep=2.2pt, line width=0.8pt},
  fase/.style={font=\tiny\itshape, text=gray, rotate=90, anchor=south}]
\node[box] (dados) {\textbf{Dados} UCI 697: 3.630 alunos com desfecho (evadiu $\times$ formou)};
\node[box, below=of dados] (prep) {\textbf{Janelas da trajetória}: matrícula (17 var.) $\to$ fim do 1º sem. (26) $\to$ fim do 2º sem. (32)\\curso e macroeconomia fora; categorias com $<30$ alunos $\to$ ``outros''};
\node[box, below=of prep] (dist) {\textbf{Distância} entre alunos: Gower $|$ associação (\textit{u\_dep})};
\node[box, below=of dist] (grafo) {\textbf{Grafo} $G{=}(V,E,w)$: kNN mútuo com $k{=}40$ e peso $w{=}1{-}d/d_{\max}$\\\textit{o desfecho não entra no grafo}};
\node[meia, draw=corLeiden, below=of grafo, xshift=-0.2225\columnwidth] (lei) {\textcolor{corLeiden}{\textbf{Leiden}} (V1, V3)\\comunidades: que perfis existem?};
\node[meia, draw=corSage, below=of grafo, xshift=0.2225\columnwidth] (sage) {\textcolor{corSage}{\textbf{GraphSAGE}} (V2, V4)\\a vizinhança ajuda a prever?};
\node[box, below=of lei, xshift=0.2225\columnwidth] (aval) {\textbf{Avaliação comum}: mesmas 50 divisões 80/20 por janela\\F1 macro, AUC e revocação da evasão\\testes de Friedman e Wilcoxon entre as variações};
\node[box, below=of aval] (esp) {\textbf{Avaliação específica}\\\textcolor{corLeiden}{Leiden}: $Q$, V de Cramér, permutação, homofilia\\\textcolor{corSage}{GraphSAGE}: estabilidade entre divisões, distâncias e valores de $k$};
\node[box, below=of esp] (interp) {\textbf{Interpretação}: perfis, robustez, pares com destinos opostos};
\foreach \a/\b in {dados/prep, prep/dist, dist/grafo, lei/aval, sage/aval, aval/esp, esp/interp}{\draw[->] (\a) -- (\b);}
\draw[->] (grafo.south) ++(-0.2225\columnwidth,0) -- (lei.north);
\draw[->] (grafo.south) ++(0.2225\columnwidth,0) -- (sage.north);
\node[fase] at ([xshift=-1mm]prep.west) {preparação};
\node[fase] at ([xshift=-1mm, yshift=-2mm]grafo.west) {modelagem};
\node[fase] at ([xshift=-1mm, yshift=-3mm]aval.west) {avaliação};
\end{tikzpicture}
\caption{Fluxo do protótipo nas fases da CRISP-NET. O desfecho só é usado na avaliação e na interpretação. As cores dos métodos são as mesmas em todas as figuras.}
\label{fig-fluxo}
\end{figure}

\begin{table}[!ht]
\centering
\caption{As quatro variações: distância usada para construir o grafo (linhas) e método aplicado a ele (colunas).}
\label{tab-variacoes}
\small
\begin{tabular}{@{}lcc@{}}
\toprule
 & \textbf{\textcolor{corLeiden}{Leiden}} & \textbf{\textcolor{corSage}{GraphSAGE}} \\
 & \footnotesize que perfis existem? & \footnotesize a vizinhança ajuda a prever? \\
\midrule
\textbf{Gower}      & variação 1 (V1) & variação 2 (V2) \\
\textbf{Associação} & variação 3 (V3) & variação 4 (V4) \\
\bottomrule
\end{tabular}
\end{table}

\subsection{Modelagem: dados e grafo}

\textbf{Dados.} Os 4.424 registros têm 36 atributos, sem valores ausentes nem duplicados. Usamos só as classes com desfecho final: 3.630 alunos, dos quais 1.421 (39,1\%) evadiram e 2.209 se formaram. Os 794 \textit{Enrolled} ficam fora porque ainda estão no curso: tratá-los como classe misturaria ``desfecho desconhecido'' com desfecho. O alvo é 1 (evadiu) ou 0 (formou).

\textbf{Janelas da trajetória.} O conjunto registra o aluno em momentos sucessivos: a matrícula, o 1º semestre, o 2º semestre e, por fim, o desfecho. Cada janela usa só o que a instituição já sabia naquele momento, e o desfecho é sempre o alvo. A \textit{matrícula} tem 13 variáveis categóricas (estado civil, modo de ingresso, turno, qualificação anterior, nacionalidade, escolaridade e ocupação da mãe e do pai, deslocado, necessidades especiais, gênero, internacional) e 4 numéricas (ordem de candidatura, nota da qualificação anterior, nota de admissão, idade). O \textit{fim do 1º semestre} acrescenta devedor, mensalidade em dia, bolsista e as 6 medidas do semestre (unidades creditadas, inscritas, avaliadas, aprovadas, sem avaliação e nota). O \textit{fim do 2º semestre} acrescenta as mesmas 6 medidas do 2º semestre. Quanto mais tarde a janela, mais informação e menos tempo para agir; por isso comparamos as três. Há uma ressalva para a última: 51\% dos evadidos já têm zero aprovações no 2º semestre (contra 3\% dos formados), ou seja, parte dessa janela registra a evasão já em curso.

Duas exclusões são justificadas pelos dados. Os indicadores macroeconômicos têm só 10 combinações distintas, uma por coorte de ingresso, e descrevem o ano, não o aluno. O curso é contexto institucional, não característica do aluno; mantê-lo na distância agruparia os alunos por curso antes de agrupá-los por perfil. Categorias com menos de 30 alunos (0,8\% a 2,8\% dos alunos por variável) são agrupadas em ``outros'', para que categorias raríssimas não dominem a distância.

\textbf{Distâncias.} A de Gower (Equação~\ref{eq-gower}) é a referência para dados mistos. A distância por associação segue o \textit{preset} \texttt{u\_dep} do \texttt{manydist}: cada categoria é descrita pela distribuição dos alunos que a têm nas demais variáveis, e a distância entre duas categorias é a variação total entre essas distribuições; as numéricas são padronizadas e rotacionadas por PCA; cada termo é dividido pela sua distância média. O grupo reimplementou o \textit{preset} em Python e o conferiu contra o pacote R (correlação de 1,000 em 500 alunos).

\textbf{Grafo.} O modelo é um grafo simples, não direcionado, ponderado e homogêneo $G = (V, E, w)$, um por janela e distância. \textit{Vértices}: cada vértice é um aluno, a unidade sobre a qual a instituição decide. \textit{Arestas}: $\{i, j\} \in E$ se $j$ está entre os $k$ alunos mais próximos de $i$ \textit{e} $i$ entre os $k$ mais próximos de $j$ (kNN mútuo), pela distância $d$ calculada só com as variáveis da janela. Como não há vínculo real entre alunos, a relação modelada é a semelhança; o desfecho fica fora da construção, o que evita vazamento de informação. A versão mútua torna a relação simétrica, liga só pares de semelhança recíproca e impede vértices concentradores (o grau máximo é $k$). \textit{Pesos}: $w(i,j) = 1 - d(i,j)/d_{\max} \in (0, 1]$; a divisão por $d_{\max}$ põe as duas distâncias na mesma escala. \textit{Atributos dos vértices}: as variáveis da janela (categóricas em \textit{one-hot}, numéricas padronizadas). \textit{Escolha de $k = 40$}: $k$ controla a conectividade. Valores da ordem de $\ln N \approx 8$ fragmentam o grafo, e valores perto de $\sqrt{N} \approx 60$ ligam alunos pouco parecidos. Com $k = 20$, até 80 alunos ficam isolados e o grafo se parte em até 92 componentes; com $k = 60$, a modularidade cai e as comunidades se fundem. Com $k = 40$ (1,1\% de $N$), 95\% a 99,9\% dos alunos ficam no maior componente (Tabela~\ref{tab-grafo}); a Seção~\ref{ssec:robustez} mostra que os resultados se mantêm com $k = 20$ e $k = 60$.

\subsection{Método de solução}

\textbf{Leiden (V1 e V3).} Usamos o Leiden com partição por modularidade ponderada (Equação~\ref{eq-modularidade}; biblioteca \texttt{leidenalg}), por garantir comunidades conectadas \citep{traag2019}. Ele encontra as comunidades sem olhar o desfecho e responde ``que perfis existem''. Para entrar na avaliação comum, cada aluno de teste recebe como probabilidade de evadir a taxa de evasão da sua comunidade, calculada só com os alunos de treino e suavizada para a taxa geral (como se cada comunidade tivesse mais 5 alunos ``médios''), para que comunidades pequenas não deem 0\% ou 100\% por acaso.

\textbf{GraphSAGE (V2 e V4).} A GraphSAGE (Equação~\ref{eq-sage}) responde ``a vizinhança ajuda a prever''. Foi escolhida por ser indutiva: pode classificar um ingressante novo sem retreino, condição para uso institucional. Neste estudo ela é usada no modo transdutivo, com o grafo de todos os alunos (ver a avaliação comum); o uso indutivo fica para a aplicação. Usamos 2 camadas de 64 neurônios, \textit{dropout} de 0,5, Adam (taxa 0,01, \textit{weight decay} $5\cdot10^{-4}$), até 300 épocas com parada antecipada (paciência de 30) pelo F1 em 15\% do treino, e pesos de classe inversos à frequência \citep{fey2019}. A configuração é padrão e não foi ajustada.

\textbf{Avaliação em duas camadas.} Leiden e GraphSAGE têm naturezas diferentes (um agrupa, o outro classifica). Por isso, a avaliação tem uma camada \textit{comum}, que coloca as quatro variações na mesma régua, e uma camada \textit{específica}, que mede o que cada método promete.

\textit{(a) Avaliação comum.} As quatro variações são testadas nas \textit{mesmas} divisões: validação cruzada estratificada de 5 dobras (cada dobra separa 80\% para treino e 20\% para teste), repetida 10 vezes, totalizando 50 divisões 80/20 por janela. Repetir evita que um sorteio favorável decida a comparação \citep{shchur2018}. Cada método aprende só com os 80\% de treino; o grafo usa todos os alunos, mas nenhum desfecho (modo transdutivo, padrão em classificação de vértices). A regra de decisão é balanceada nos dois métodos: no Leiden, o limiar é a taxa de evasão do treino; na GraphSAGE, usam-se pesos de classe. As métricas são o F1 macro (não favorece a classe majoritária), a AUC (qualidade da ordenação do risco, sem depender do limiar) e a revocação da evasão (fração dos evadidos encontrados). Como piso, um classificador que sempre prevê ``formou'' tem F1 de 0,38 e AUC de 0,50. As diferenças são testadas pelo teste de Friedman (quatro variações) e pelo de Wilcoxon pareado por divisão, entre pares de variações que diferem em um só fator (distância ou método); como as divisões compartilham alunos, os valores-$p$ são otimistas, e a decisão também considera o tamanho da diferença e a fração de divisões em que ela se repete.

\textit{(b) Avaliação específica.} Para o Leiden: modularidade $Q$; V de Cramér entre comunidade e desfecho,
\begin{equation}
V = \sqrt{\chi^2 / \big(n\,(\min(r,c) - 1)\big)},
\label{eq-cramer}
\end{equation}
nas comunidades com pelo menos 30 alunos, com teste de permutação (1.000 embaralhamentos do desfecho com o grafo fixo); e homofilia (fração de arestas entre alunos com o mesmo desfecho) e assortatividade. Para a GraphSAGE, que usa as arestas na própria agregação (Equação~\ref{eq-sage}): a estabilidade do desempenho entre as 50 divisões (desvio-padrão), entre as duas distâncias e entre valores de $k$ (Seção~\ref{ssec:robustez}). Para interpretar, descrevemos cada comunidade pelas variáveis que mais se afastam da média, comparamos as partições das duas distâncias pelo índice de Rand ajustado (ARI) e analisamos pares de vizinhos com desfechos opostos.

\textbf{Protótipo.} Escrito em Python (igraph, leidenalg, PyTorch Geometric, scikit-learn), recebe o arquivo do conjunto de dados e gera todas as tabelas e figuras. O núcleo (dados, janelas, distâncias, grafo, Leiden, GraphSAGE) fica em \texttt{src/grafo\_evasao.py}; \texttt{avaliacao\_unica.py} executa a avaliação comum; \texttt{metricas\_leiden.py}, a específica; \texttt{resumo\_avaliacao.py} faz os testes e as figuras, com a paleta única de \texttt{src/cores.py}. O kNN desempata pelo índice do aluno e o Leiden usa semente fixa, o que reproduz exatamente as partições; as redes neurais usam semente por divisão.
